\vspace{-4pt}
\section{Related Work}\label{sec:related}

Prior coding benchmarks establish reliable end state scoring for code changes, but they offer limited evidence about the loops that sustain long horizon execution. Function level code generation benchmarks~\citep{chen2021evaluating} pioneered programmatic evaluation on localized snippets, while repository level benchmarks centered on SWE-bench~\citep{jimenez2023swe} package work as issue resolution and emphasize end state pass rate.
Feature and evolution benchmarks~\citep{zhou2026featurebench} enlarge task scope, and broader long horizon benchmarks~\citep{kwa2025measuring} extend duration, repository scale, or interaction depth.
Across these families, prerequisite state, residual work routing, and regression obligations generally remain outside the executable scoring contract. Success is therefore assigned to submitted artifacts rather than to the loop state that produced them.
Modern coding agent infrastructure~\citep{claude_code,openai_codex_prompting,anthropic_claude_code_workflows,huntley_ralph_loop} increasingly relies on persistent objectives, worker orchestration, restarts, and context management. Loop implementation therefore becomes part of the evaluated behavior, while current benchmark designs provide limited evidence about how evaluated loops preserve state and route residual work.
Software engineering work on change impact and release planning~\citep{lehnert2011taxonomy} treats dependency structure as an explicit artifact. \lhb{} draws on this view by converting dependency structure into executable obligations for coding agent loop evaluation. Appendix~\ref{app:related-bench} provides the full grouping.
